\documentclass[12pt]{article}
\usepackage[utf8]{inputenc}
\usepackage[T1]{fontenc}
\usepackage{hyperref}
\title{T\`ay Tac Language Resources for K-12 Curriculum Development}
\author{Compiled by Claude Sonnet 5 (claude-sonnet-5) for Andr\'as Kornai}
\date{}
\begin{document}
\maketitle

\section{Introduction}
T\`ay Tac (White Tay; ISO 639-3 tyt, Glottocode tayt1241) is a Southwestern Tai language spoken in M\uhorn{}\`ong T\'ac district in eastern S\ohorn{}n La Province, Vietnam. Search results say the population is unknown and may be very small, and that the group may have merged into neighboring groups. They also say the Tay Tac are recognized as White Tay, with a language and script of the Tay Daeng type. The searches found very little.

\section{Primary Educational Resources}
No primer, textbook or school programme turned up.

\section{Linguistic and Reference Materials}
\begin{itemize}
\item T\`ay Tac language (Wikipedia). It lists Latin and Tai Viet scripts: \url{https://en.wikipedia.org/wiki/T%C3%A0y_Tac_language}
\item Ferlus, a French-language study of Tai Ngh\^e An, which mentions the Tay Tac: \url{https://halshs.archives-ouvertes.fr/halshs-01617791v1/file/Ferlus2008TaiNgh%C3%AAAn_version2017.pdf}
\item T\`ay language (Wikipedia), for the wider T\`ay context. This is a different, though related, language: \url{https://en.wikipedia.org/wiki/T%C3%A0y_language}
\end{itemize}

\section{Religious and Community Resources}
\begin{itemize}
\item Joshua Project profile of the Tay Tac in Vietnam: \url{https://joshuaproject.net/people_groups/15308}
\end{itemize}

\section{Digital and Audio Resources}
None found. Some ASJP source pages appeared in the results, but I did not confirm that they cover this language.

\section{Curriculum Development Context}
The documentation is thin, the number of speakers is unclear, and no school use was found. The first step would be to establish with the community whether the language is still spoken. The related Tay Daeng and Tai Viet materials may help, but should not be treated as material for Tay Tac.
\end{document}